% Chapter 16: Computational vs. Information-Theoretic Learnability
% Centerpiece: the Kearns-Valiant 1994 gap — classes with finite VC dim
% that are NOT efficiently learnable (under cryptographic assumptions).
% Concepts: computational_hardness, proper_improper_separation, requires_assumption edges.
% SHORT chapter (~3-4 pages).

\chapter{Computational vs.\ Information-Theoretic Learnability}
\label{ch:computational}

The Fundamental Theorem of \Cref{ch:pac} characterizes learnability in
information-theoretic terms: $\mathcal{H}$ is PAC learnable if and only if
$\VC(\mathcal{H}) < \infty$.  The characterization is silent about
\emph{computation}.  It guarantees the \emph{existence} of a sample size
$m(\varepsilon, \delta)$ and a learner $A$ that achieves accuracy $\varepsilon$
with confidence $1 - \delta$, but it says nothing about how long $A$ takes to
run.

This silence conceals a gap.  There exist hypothesis classes with finite VC
dimension, hence information-theoretically learnable, that no polynomial-time
algorithm can PAC learn, unless widely believed cryptographic assumptions are
false.  The gap between what is \emph{learnable} and what is \emph{efficiently
learnable} is the subject of this chapter.  It is a negative result, but a
structurally important one: it shows that information-theoretic
characterizations, however elegant, do not tell the whole story.

\medskip

The chapter has three sections.

\begin{enumerate}[leftmargin=*,itemsep=2pt]
\item \textbf{The information--computation gap} (\Cref{sec:gap}): the
  Kearns--Valiant result and the cryptographic assumptions it requires.
\item \textbf{Proper vs.\ improper learning} (\Cref{sec:proper-improper-comput}): the
  separation between proper and improper learnability, and why the
  representation matters.
\item \textbf{The \rel{requires\_assumption} edges}
  (\Cref{sec:requires-assumption}): what this edge type means structurally in
  the concept graph.
\end{enumerate}


% ================================================================
% SECTION 1: THE INFORMATION--COMPUTATION GAP
% ================================================================

\section{The Information--Computation Gap}
\label{sec:gap}

\begin{graphnode}[title={Graph Node: \nde{computational\_hardness}}]
Layer: \texttt{meta}.  Incoming edge:
$\edge{computational\_hardness}{requires\_assumption}{cryptographic\_hardness}$.
This node represents the structural fact that efficient learnability requires
assumptions beyond finite VC dimension, specifically, cryptographic hardness.
\end{graphnode}

The central result is due to Kearns and Valiant~\cite{KearnsValiant1994}.

\begin{thm}[Kearns--Valiant 1994]
\label{thm:kearns-valiant}
Under a cryptographic hardness assumption (the original construction uses the
hardness of factoring; later strengthenings use decisional composite
residuosity or learning with errors, detailed below), there exist concept
classes $\mathcal{C}$ with:
\begin{enumerate}[leftmargin=*,nosep]
\item $\VC(\mathcal{C}) < \infty$, so that $\mathcal{C}$ is PAC learnable
  information-theoretically,\formal{FLT_Proofs.Theorem.PAC.fundamental_theorem} but
\item no polynomial-time algorithm PAC learns $\mathcal{C}$.
\end{enumerate}
\end{thm}

The original construction uses Boolean formulae: the class of polynomial-size
Boolean formulae has VC dimension polynomial in the number of variables, yet
PAC learning this class is at least as hard as breaking certain public-key
cryptosystems.  The reduction goes through the construction of
\emph{cryptographic pseudorandom function families}: if a learner could
efficiently distinguish the target concept from random, it could break the
cryptographic assumption.

\begin{remark}[The structure of the reduction]
\label{rmk:kv-reduction}
The Kearns--Valiant reduction is from a cryptographic primitive to a learning
problem, not from a worst-case problem.  The learner is given random examples
$(x, c(x))$ where $x$ is drawn from a distribution and $c$ is the target
concept.  If the learner succeeds in learning, it implicitly solves the
cryptographic problem on random instances.  This is why the hardness is
\emph{average-case}, not worst-case.
\end{remark}

\subsection{Why P $\neq$ NP Does Not Suffice}

One might hope that the information--computation gap follows from
$\mathrm{P} \neq \mathrm{NP}$.  It does not.  Applebaum, Barak, and
Xiao~\cite{ApplebaumBarakXiao2008} showed the following.

\begin{thm}[Applebaum--Barak--Xiao 2008]
\label{thm:abx}
There is no \emph{black-box} reduction from $\mathrm{NP}$-hardness to the
hardness of PAC learning.  Specifically, if PAC learning a class $\mathcal{C}$
is hard, this hardness cannot be established by a reduction that treats the
learner as an oracle, unless $\mathrm{NP} \subseteq \mathrm{coAM}$.
\end{thm}

The reason is structural.  PAC learning is an \emph{average-case} problem: the
learner receives random examples from a distribution and must generalize.
$\mathrm{NP}$-hardness is a \emph{worst-case} notion.  Reductions from
worst-case to average-case problems are notoriously difficult and, in many
settings, provably impossible via relativizing techniques.

\subsection{The Necessary Assumptions}

The assumptions used in the literature fall into three families, in
roughly increasing order of generality:

\begin{enumerate}[leftmargin=*,itemsep=3pt]
\item \textbf{Factoring and RSA.}  The original Kearns--Valiant result
  assumes that factoring $n$-bit integers requires super-polynomial time.
  This is the most classical assumption but also the most fragile: it is
  broken by quantum computers (Shor's algorithm).

\item \textbf{Decisional Composite Residuosity (DCRA).}  Used in subsequent
  refinements, DCRA assumes that it is hard to distinguish $n$th residues
  modulo $n^2$.  Like factoring, this is vulnerable to quantum attacks.

\item \textbf{Learning With Errors (LWE).}  Introduced by
  Regev~\cite{Regev2005}, the LWE assumption states that it is hard to recover
  a secret vector $\mathbf{s}$ from noisy inner products
  $\langle \mathbf{a}_i, \mathbf{s} \rangle + e_i \pmod{q}$.  LWE is
  believed to resist quantum computers and has become the standard assumption
  in modern hardness-of-learning results~\cite{Daniely2021}.
\end{enumerate}

\begin{obstbox}
The information--computation gap is an \emph{obstruction} in the concept
graph: the edge
$\edge{vc\_dimension}{characterizes}{pac\_learnability}$ is
information-theoretically valid, but no analogous computational edge is known.
Whether \emph{any} combinatorial dimension characterizes efficient PAC
learnability is open; no such dimension has been found, and no proof rules
one out.  What \emph{is} settled, and machine-checked, is a sharper local
fact: finite VC dimension does not control the statistical-query dimension.
Singleton indicators on $\N$ have $\VC = 1$ yet infinite SQ
dimension,\formal{FLT_Proofs.Theorem.Extended.vcdim_not_implies_hardness} so
VC dimension already fails to capture hardness in the SQ model, one concrete
instance of the gap, even though its full extent is unknown.
\end{obstbox}


% ================================================================
% SECTION 2: PROPER VS. IMPROPER LEARNING
% ================================================================

\section{Proper vs.\ Improper Learning}
\label{sec:proper-improper-comput}

The information--computation gap has a subtler cousin: the gap between
\emph{proper} and \emph{improper} learning.

\begin{defn}[Proper and Improper Learning]
\label{def:proper-improper}
A PAC learner for $\mathcal{H}$ is \emph{proper} if it always outputs a
hypothesis $h \in \mathcal{H}$.  It is \emph{improper} if it may output any
efficiently evaluable function $h$, not necessarily in $\mathcal{H}$.
\end{defn}

The distinction matters because the \emph{representation} of hypotheses affects
computational complexity.

\begin{thm}[Proper--Improper Separation]
\label{thm:proper-improper}
There exist concept classes $\mathcal{C}$ such that:
\begin{enumerate}[leftmargin=*,nosep]
\item proper PAC learning of $\mathcal{C}$ is $\mathrm{NP}$-hard (so, unless
  $\mathrm{RP} = \mathrm{NP}$, no efficient proper learner exists), but
\item there exists a polynomial-time improper learner for $\mathcal{C}$.
\end{enumerate}
\end{thm}

The canonical example is the class of $3$-term DNF formulae.  Finding a
consistent $3$-term DNF is $\mathrm{NP}$-hard~\cite{PittValiant1988}, but the
class of $3$-term DNFs is contained in the class of $3$-CNF formulae, which
\emph{is} efficiently learnable.  An improper learner simply outputs a $3$-CNF
formula that is consistent with the data.  The hypothesis is not a DNF, but it
classifies correctly.

Information-theoretically, a proper learner always exists: any class of finite
VC dimension is learned by proper empirical risk minimization, which returns a
consistent member of the class whenever one
exists.\formal{FLT_Proofs.Complexity.Compression.vcdim_finite_imp_proper_finite_support_learner}
The separation is therefore purely computational: the proper learner exists
but cannot be run efficiently, whereas the improper learner can.

\begin{remark}[Representation vs.\ class]
\label{rmk:representation}
The proper--improper separation illustrates a principle that runs throughout
computational learning theory: what matters for computational complexity is not
the \emph{set of functions} $\mathcal{H} \subseteq \{0,1\}^X$ but the
\emph{representation} of those functions.  Two representations of the same set
of Boolean functions can have dramatically different computational properties.
This is why the VC dimension, which depends only on the set
$\mathcal{H}$, cannot capture computational learnability.
\end{remark}

The structural consequence is that efficient learnability is a property of
\emph{represented} hypothesis classes, not of hypothesis classes as sets.  The
concept graph captures this through the node \nde{proper\_improper\_separation},
which has type \texttt{separation} and connects to
\nde{computational\_hardness} via a \rel{motivates} edge.


\subsection{The Transformer Face of the Gap}
\label{sec:attention-gap}

The information--computation gap is not a museum piece; it is the daily
condition of deep learning.  A softmax-attention head induces a concept class
that is machine-checked to be well-behaved, hence of finite VC dimension and
information-theoretically PAC
learnable.\formal{TLT.Bridge.softAttention_wellBehaved} A proper learner for
it therefore exists in principle.  Whether that learner runs in polynomial
time is the open, and widely doubted, question: hardness-of-learning results
for neural networks under LWE~\cite{Daniely2021} place architectures of this
kind on the conjecturally efficient-unlearnable side of
\Cref{thm:kearns-valiant}.

The proper/improper distinction is what makes practice possible. Recovering
the exact parameters of a target attention head, which is proper learning, is the
regime where hardness bites. Gradient descent does something improper
instead: it returns \emph{a} network that predicts well, not the target
network, exactly as the $3$-CNF learner returns a CNF rather than the target
DNF. The gap of this chapter is, in this sense, the precise reason that
training large models is an improper-learning enterprise.


% ================================================================
% SECTION 3: THE REQUIRES_ASSUMPTION EDGES
% ================================================================

\section{The \rel{requires\_assumption} Edges}
\label{sec:requires-assumption}

The concept graph contains several edges of type \rel{requires\_assumption}.
These edges encode a specific logical structure: the source node's validity
depends on an assumption that is not provable from the axioms of the theory
alone.

\begin{defn}[The \rel{requires\_assumption} edge type]
\label{def:requires-assumption}
An edge $\edge{A}{requires\_assumption}{B}$ means: the result or property
associated with node $A$ holds \emph{conditional on} the assumption encoded
by node $B$.  If $B$ is falsified, the status of $A$ becomes indeterminate.
\end{defn}

The principal instances in the graph are:

\begin{enumerate}[leftmargin=*,itemsep=3pt]
\item $\edge{computational\_hardness}{requires\_assumption}{cryptographic\_hardness}$.
  The information--computation gap of \Cref{thm:kearns-valiant} holds only
  under cryptographic assumptions (hardness of factoring, or LWE).  Without
  them it remains conceivable, though it would be a startling
  collapse, that every class of finite VC dimension is efficiently learnable.
  The assumption is genuinely cryptographic: the hardness of a specific
  computational problem, not a constraint on the learner's hypothesis space.

\item $\edge{posterior\_consistency}{requires\_assumption}{bayesian\_inference}$.
  Posterior consistency (\Cref{ch:agents}, \Cref{sec:bayesian-pac}) holds under
  conditions on the prior and the model class.  When these conditions
  fail, for example when the prior assigns zero mass to neighborhoods of
  the true parameter, posterior consistency fails
  (Diaconis--Freedman~\cite{DiaconisFreedman1986}).  The edge records this
  conditional dependence.
\end{enumerate}

\begin{remark}[Structural meaning]
\label{rmk:requires-assumption-structure}
The \rel{requires\_assumption} edge type is the graph's mechanism for
representing \emph{conditional} theorems.  Unlike \rel{implies} edges, which
are unconditional, a \rel{requires\_assumption} edge signals that the source
node's content is only as secure as the target node's assumption.  This is
the graph-theoretic encoding of the phrase ``under the assumption that\ldots''
that pervades computational learning theory.
\end{remark}


% ================================================================
% OPEN PROBLEMS
% ================================================================

\section{Open Problems}
\label{sec:computational-open}

The computational layer is where formal learning theory is least complete:
its central hardness results are conditional, and its central
characterization question is open.  Each problem below is tagged a
\emph{known unknown} (KU), sharply posed but unresolved, or an \emph{unknown
known} (UK), where the right object is still missing; a design-lab or
literature anchor is named where one exists.

\begin{openprob}[A dimension for efficient learnability]
\label{op:efficient-dimension}
Is there \emph{any} combinatorial parameter whose finiteness characterizes
efficient (polynomial-time) PAC learnability, as VC dimension characterizes
the information-theoretic case?  None is known, and none is ruled out.  A
known unknown, the central question of the computational layer.
\end{openprob}

\begin{openprob}[Formalizing the cryptographic gap]
\label{op:formalize-kv}
The information-theoretic side of \Cref{thm:kearns-valiant} is machine-checked
(\leanname{vcdim_finite_imp_proper_finite_support_learner}); the cryptographic
hardness side is not.  Can a Kearns--Valiant-style hardness result, a finite-VC
class with no efficient learner under LWE, be formalized in a proof
assistant?  An unknown known~\cite{KearnsValiant1994,Daniely2021}.
\end{openprob}

\begin{openprob}[Distribution-free SQ hardness, verified]
\label{op:sq-distribution-free}
The machine-checked separation finite-VC$\,\not\Rightarrow\,$finite-SQ uses a
fixed witness distribution (\leanname{vcdim_not_implies_hardness}).  Can one
verify a finite-VC class whose SQ dimension is super-polynomial under
\emph{every} distribution, matching the parities lower
bound~\cite{BlumFurstJacksonKearnsMansourRudich1994}?  A known unknown.
\end{openprob}

\begin{openprob}[Is attention efficiently learnable?]
\label{op:attention-efficient}
A softmax-attention class is information-theoretically learnable
(\leanname{softAttention_wellBehaved}).  Is it efficiently PAC learnable, or
does learning attention reduce to a cryptographic problem?  A known unknown on
the transformer side, and the modern instance of \Cref{thm:kearns-valiant}.
\end{openprob}

\begin{openprob}[Past the black-box barrier]
\label{op:worst-to-average}
Applebaum--Barak--Xiao rule out \emph{black-box} reductions from
$\mathrm{NP}$-hardness to learning hardness.  Is there a non-black-box route
that bases hardness of learning on a worst-case assumption such as
$\mathrm{P} \neq \mathrm{NP}$?  A known unknown~\cite{ApplebaumBarakXiao2008}.
\end{openprob}

\begin{openprob}[Characterizing improper learnability]
\label{op:improper-characterization}
Efficient \emph{improper} learning can be strictly easier than proper learning
(\Cref{thm:proper-improper}).  Is there a structural characterization of the
classes that are efficiently improperly learnable?  An unknown known: no
combinatorial handle on improper efficiency is known.
\end{openprob}

\begin{openprob}[Minimal assumptions for hardness]
\label{op:minimal-assumption}
Hardness-of-learning currently rests on strong primitives (factoring, LWE).
Is some cryptographic assumption \emph{necessary} for the
information--computation gap: does an efficient learner for every finite-VC
class imply the collapse of one-way functions?  A known
unknown~\cite{Daniely2021}.
\end{openprob}

\begin{openprob}[Learning implies cryptography?]
\label{op:learning-to-crypto}
The Kearns--Valiant reduction builds hardness \emph{from} cryptography. Does
the converse hold: does the hardness of learning a natural concept class
\emph{yield} a cryptographic primitive?  An unknown known: a generic
learning-to-cryptography transformation is not established.
\end{openprob}

\begin{openprob}[Size of the proper/improper gap]
\label{op:proper-improper-size}
For \Cref{thm:proper-improper}, how large can the separation be: how much
harder is proper than improper learning across natural classes, measured in
running time?  A known unknown~\cite{PittValiant1988}.
\end{openprob}

\begin{openprob}[A proper/improper gap for attention]
\label{op:attention-proper-improper}
Does an attention class exhibit a proper/improper gap: is there an efficient
improper learner (a surrogate model) even where recovering the target head is
intractable, formalizing the intuition that gradient descent learns
improperly?  A known unknown on the transformer side
(\leanname{softAttention_wellBehaved},
\leanname{vcdim_finite_imp_proper_finite_support_learner}).
\end{openprob}

\bigskip

This chapter has established the structural gap between
information-theoretic and computational learnability.  The gap is real: it
produces classes that are provably learnable but (conditionally) not
efficiently learnable.  And no combinatorial dimension is yet known to
characterize efficient PAC learnability; whether one can exist is itself an
open problem.  The next chapter examines a different kind of gap: the distance
between binary classification and the richer learning settings that require
new mathematical primitives.
